% !PRINTED 338
\weberpage{338}

\begin{center}
\gesperrt{Eleventh Chapter.}

\textbf{\Large Ideals in Quadratic Fields.}

\rule{3cm}{0.4pt}
\end{center}

\begin{center}
\textbf{\S\ 90. \gesperrt{Discriminant of the Quadratic Field.}}
\end{center}

A quadratic field arises when a square root $\sqrt{d}$ is \gesperrt{adjoined} to the
field of rational numbers; here $d$ may be assumed to be \gesperrt{an integer without
a square divisor}. The field is real or imaginary according as $d$ is positive or negative. If we denote the field of
rational numbers, the absolute domain of rationality, by $\mathfrak{R}$, and
by $\mathfrak{R}(x,x',\dots)$ the field which arises from $\mathfrak{R}$ by the adjunction of $x,x',\dots$,
then the quadratic field $\Omega$ is to be denoted thus:
\weq{1}{\Omega=\mathfrak{R}(\sqrt{d}).}

Every number of the field $\Omega$ can be put into the form:
\weq{2}{\omega=\frac{x+y\sqrt{d}}{2},}
in which $x, y$ are rational integral or fractional numbers.

The number \gesperrt{conjugate} to $\omega$
\weq{3}{\omega'=\frac{x-y\sqrt{d}}{2}}
is likewise contained in the field $\Omega$, and consequently $\Omega$ is a
\gesperrt{normal field.}

In order that $\omega$ be \gesperrt{an integer}, it is necessary that
\[
\omega+\omega'=x,\quad(\omega-\omega')^2=y^2 d
\]
be rational integers, and since $d$ is to have no square divisor,
$x$ and $y$ must be integers. But in order that $\omega$ should really be an
integer, the norm $\frac{1}{4}(x^2-y^2 d)$ must also be a rational integer, i.\,e.
\weq{4}{x^2-y^2 d\equiv 0\pmod{4}}
must hold. If $d\equiv 2$ or $\equiv 3\pmod 4$, this condition can only be
satisfied if $x$ and $y$ are even numbers, and hence, if we replace
$x, y$ by $2x, 2y$, it follows:
%CONT>>
